Tato sekce shrnuje praktické rubriky a postupy pro transparentní hodnocení v prostředí s možností použití generativní AI. Cílem je hodnotit nejen konečný produkt, ale i proces, ověřování výsledků a reflexi použití AI.

Principy a požadavky
\begin{itemize}
  \item Každé zadání s možností použití AI musí obsahovat pole „Deklarace použití AI“ s povinnými položkami (viz níže).
  \item Hodnocení zahrnuje procesní artefakty (milníky, drafty, logy) vedle výsledku; u vybraných výstupů požadujte stručnou verifikaci.
  \item Formativní podpora nástroji (např. Copilot) může učiteli šetřit čas, vždy však vyžadujte lidské přezkoumání \cite{kb_ai_101_tipu_pdf_04e4a115_page_21_2}.
\end{itemize}

Povinné pole pro deklaraci AI (šablona)
\begin{enumerate}
  \item Použité nástroje a verze.
  \item Prompt nebo export konverzace.
  \item Co bylo převzato a co upraveno.
  \item Jak bylo ověřeno (stručně).
  \item Reflexe: 2–3 věty o vlivu AI na postup a výsledek.
\end{enumerate}

Rubrika — klíčové položky
\begin{description}
  \item[Proces a dokumentace] Kvalita milníků, draftů a logů; sledovatelnost postupu.
  \item[Deklarace AI a verifikace] Kompletní deklarace + stručné doklady ověření.
  \item[Technická kvalita] Správnost řešení, argumentace, metodika.
  \item[Reflexe a etika] Sebereflexe role AI, etické úvahy.
  \item[Originalita a aplikace] Originalita a lokální uplatnění poznatků.
\end{description}

Návrh vah (ilustrační)
\begin{itemize}
  \item Proces a dokumentace: 25\,\%
  \item Deklarace AI a verifikace: 10\,\%
  \item Technická kvalita: 45\,\%
  \item Reflexe a etika: 10\,\%
  \item Originalita a aplikace: 10\,\%
\end{itemize}

Milestone‑based hodnocení — implementační vzor
\begin{enumerate}
  \item Definujte 3–5 milníků s konkrétními výstupy (návrh, prototyp, testy, závěrečná zpráva + reflexe).
  \item U každého milníku určete váhu a požadované artefakty (draft, log, export konverzace).
  \item Požadujte pro každý milník „Deklaraci AI“ a krátkou verifikaci (testy, zdroje, porovnání).
  \item Zaveďte formativní kontrolní body s komentáři učitele/TA a provádějte náhodné spot‑checky (např. 5–10\,\%).
\end{enumerate}

Příklad rozložení milníků (ilustrativně)
\begin{itemize}
  \item A: Návrh a plán verifikace (10\,\%).
  \item B: Prototyp a logy použití (30\,\%).
  \item C: Testy a verifikace výsledků (30\,\%).
  \item D: Závěrečná zpráva + deklarace AI + reflexe (30\,\%).
\end{itemize}

Praktická šablona pro formativní zpětnou vazbu (kopírovatelný prompt)
\begin{verbatim}
Jsi asistent pro zpětnou vazbu. Analyzuj studentský postup a:
1) vypiš max. 3 konkrétní chyby nebo nejasnosti;
2) pro každou chybu navrhni krátké cvičení nebo kontrolní otázku (max. 2 věty);
3) pokud chybí krok, požádej o doplnění.
Označ tagy: [Chyba], [Cvičení], [Požadavek].
\end{verbatim}

Kontrolní seznam pro zavedení rubriky
\begin{enumerate}
  \item Do sylabu přidejte pole „Deklarace AI“ a očekávání ke každému milníku.
  \item Připravte rubriku pro proces, deklaraci AI a verifikaci.
  \item Pilotujte postup na malé skupině a sbírejte zpětnou vazbu od studentů i TA.
  \item Stanovte mechanismus spot‑checků a eskalace podezření na neférové použití.
  \item Dokumentujte pilot a upravte rubriku před plošným nasazením.
\end{enumerate}

Krátké doporučení k transparentnosti
\begin{itemize}
  \item Studentům jasně sdělte, kde je AI povolena a jak bude její použití hodnoceno.
  \item Povinné deklarace a dokumentace promptů/konverzací zvyšují ověřitelnost a podporují zodpovědné používání nástrojů.
\end{itemize}

Poznámka: Hromadné generování formulací formativní zpětné vazby nástroji typu Copilot může ušetřit čas; vždy však vyžadujte lidskou validaci generovaných komentářů \cite{kb_ai_101_tipu_pdf_04e4a115_page_21_2}.